\subsection{Критерий Коши существования предела функции в точке}

\textbf{Теорема (Критерий Коши).} Пусть $f(x)$ определена на проколотой окрестности $\mathring U_a$.
\[ \exists \lim_{x \to a} f(x) \Leftrightarrow \forall \varepsilon > 0\; \exists \delta_\varepsilon > 0: \forall x', x'' \in (a - \delta_\varepsilon, a + \delta_\varepsilon) \setminus \{a\}\;\; |f(x') - f(x'')| < \varepsilon \]

\begin{tcolorbox}[title=Доказательство]
	$\Rightarrow$ (Необходимость). Пусть $\exists \lim_{x \to a} f(x) = A$.
	Запишем определение предела для $\frac{\varepsilon}{2}$:
	\[ \forall \varepsilon > 0\; \exists \delta_\varepsilon > 0: \forall x \in \mathring U_{\delta_\varepsilon}(a) \;\; |f(x) - A| < \frac{\varepsilon}{2} \]
	Возьмем любые $x', x''$ из этой окрестности:
	\[ |f(x') - f(x'')| = |(f(x') - A) - (f(x'') - A)| \leq |f(x') - A| + |f(x'') - A| < \frac{\varepsilon}{2} + \frac{\varepsilon}{2} = \varepsilon \]
	
	$\Leftarrow$ (Достаточность). Пусть выполнено условие Коши.
	Возьмем произвольную последовательность $\{x_n\} \subset \mathring U_a$, сходящуюся к $a$ ($x_n \to a$).
	Так как $x_n \to a$, то начиная с некоторого номера $N$, все члены попадают в $\delta_\varepsilon$-окрестность:
	\[ \forall n, m > N \implies x_n, x_m \in (a - \delta_\varepsilon, a + \delta_\varepsilon) \]
	Тогда по условию теоремы:
	\[ |f(x_n) - f(x_m)| < \varepsilon \]
	Это означает, что числовая последовательность $\{f(x_n)\}$ — фундаментальная.
	В силу полноты $\mathbb{R}$ (критерий Коши для последовательностей), она сходится.
	Так как это верно для любой последовательности, то $\exists \lim_{x\to a} f(x)$ (в смысле Гейне).
\end{tcolorbox}

\textbf{Отрицание (Критерий расходимости).}
\[ \nexists \lim_{x \to a} f(x) \Leftrightarrow \exists \varepsilon > 0\; \forall \delta > 0\; \exists x', x'' \in \mathring U_\delta(a): \;\; |f(x') - f(x'')| \geq \varepsilon \]

\textbf{Примеры отсутствия предела:}
\begin{enumerate}
	\item $f(x) = \frac{1}{x}$ на $(0, 1)$ при $x \to 0$.
	Возьмем $\varepsilon = 1$. Для любого $\delta > 0$ выберем $N$ так, что $\frac{1}{N} < \delta$.
	Положим $x' = \frac{1}{N}$ и $x'' = \frac{1}{N+1}$. Оба числа лежат в $(0, \delta)$.
	\[ |f(x') - f(x'')| = |N - (N+1)| = |-1| = 1 \ge \varepsilon \]
	Предела не существует (конечного).

	\item $f(x) = \sin \frac{1}{x}$ на $(0, 1)$ при $x \to 0$.
	Возьмем $\varepsilon = 1$. Для любого $\delta$ найдем $N$ достаточно большое.
	\[ x' = \frac{1}{\frac{\pi}{2} + 2\pi N}, \quad x'' = \frac{1}{-\frac{\pi}{2} + 2\pi N} \]
	Тогда $f(x') = \sin(\frac{\pi}{2} + 2\pi N) = 1$, а $f(x'') = \sin(-\frac{\pi}{2} + 2\pi N) = -1$.
	\[ |f(x') - f(x'')| = |1 - (-1)| = 2 \ge \varepsilon \]
	Предела не существует.
\end{enumerate}
